\section{Introduction}
Benchmarks and evaluations play a critical role in the development of large language models. They help determine which model improvements are considered useful and set the direction of future research~\cite{dotanValueladenDisciplinaryShifts2019,hutchinsonEvaluationGapsMachine2022}. Creating a benchmark requires operationalising phenomena (abstract concepts) into concrete tasks and metrics that serve as measurable proxies for model capabilities~\cite{rajiAIEverythingWhole2021a}. As an example, the `intelligence' of LLMs is frequently debated~\cite{bubeckSparksArtificialGeneral2023, mccoyEmbersAutoregressionUnderstanding2023b}, but cannot be measured directly, making it necessary to develop proxies~\cite{hernandez-oralloMeasureAllMinds2017}. The value of a benchmark depends on whether it is a good proxy for the real-world phenomenon it intends to measure. This property is known as \textit{construct validity}: the degree to which a benchmark score provides evidence for making claims about the target phenomenon~\cite{bowmanWhatWillIt2021,rajiAIEverythingWhole2021a,cronbachConstructValidityPsychological1955}. If a benchmark has high construct validity in measuring `intelligence', then a model which does well is in some sense `intelligent', but if the construct validity is low, then a high score may be irrelevant or even misleading.

The science of evaluating large language models (LLMs) is still in its early stages, with a pressing need for shared standards and best practices~\cite{weidingerEvaluationScienceGenerative2025, bowmanWhatWillIt2021}. Certain specific issues such as reproducibility and cost have been addressed via shared implementation standards~\cite{bidermanLessonsTrenchesReproducible2024a,UK_AI_Security_Institute_Inspect_AI_Framework_2024}, and item selection methods~\cite{poloTinyBenchmarksEvaluatingLLMs2024a}, respectively. Other issues, such as the best use of statistical methods~\cite{mcintoshInadequaciesLargeLanguage2024, millerAddingErrorBars2024, luettgauHiBayESHierarchicalBayesian2025} and social responsibility~\cite{bowmanWhatWillIt2021, gebruDatasheetsDatasets2021} have also been raised. \textcite{reuelBetterBenchAssessingAI2024} aggregate best practices and provide recommendations for the whole lifecycle of a benchmark. However, identifying concrete best practices for creating benchmarks with high construct validity remains a difficult task. Benchmarks with low construct validity have real consequences, since unrecognised weak links between tasks and the underlying phenomena they claim to measure can lead to poorly supported scientific claims, misdirected research, and policy implications that are not grounded in robust evidence.

Here, we assess practices around the construct validity of LLM benchmarks through a systematic review of 445 articles from leading ML and NLP conferences. The articles were coded by experts in ML and NLP using a detailed conceptual and methodological schema that identifies useful practices in the design and interpretation of benchmarks for increasing the validity of measurements. Almost all articles have weaknesses in at least one area across phenomena, tasks, metrics, and claims. Key concepts are often poorly defined or operationalised, limiting the reliability of the conclusions they draw. We call for improved practices and reporting standards for establishing construct validity in new benchmarks, and release an operational checklist of best practice recommendations.

\begin{figure}
    \centering
    \includegraphics[width=\linewidth]{images/fig-1-review-process}
    \caption{\textbf{Systematic review process.} (A) Identification and screening from relevant proceedings. (B) In-depth review and annotation of included benchmarks. A phenomenon is operationalised via a task, scored with a metric, to support a claim about this phenomenon. (C) Synthesis of best practices.}
    \label{fig:systematic-review-process}
\end{figure}

\section{Background}\label{background}

Construct validity evaluates whether an empirical test measures the phenomenon it intends to measure~\cite{alaaMedicalLargeLanguage2025a}. Formal assessments of construct validity originate from psychological testing as a means of creating tests for phenomena which cannot be directly verified, such as personality~\cite{cronbachConstructValidityPsychological1955}.

Construct validity as an overarching concept can be assessed by considering  various features of the test design~\cite{messickTestValidityMatter1998}. At the level of phenomena, \textit{face validity} considers whether a test appears prima facie a valid representation of the phenomenon~\cite{nevoFaceValidityRevisited1985,mosierCriticalExaminationConcepts1947}. At the task level, \textit{content validity} considers whether the task content represents all important aspects of the phenomenon being measured~\cite{sireciConstructContentValidity1998}. \textit{Ecological} and \textit{predictive validity} concern the relevance of the test to real-world settings~\cite{schmucklerWhatEcologicalValidity2001}, including how it predicts future performance~\cite{alaaMedicalLargeLanguage2025a}. \textit{Convergent}, \textit{discriminant}, and \textit{criterion validity} measure whether test findings correlate with, and only with, tests for similar phenomena~\cite{campbellConvergentDiscriminantValidation1959}.

With LLMs, construct validity is key for benchmarking abstract abilities such as `reasoning.' The value of construct validity has been emphasised in previous NLP literature~\cite{bowmanWhatWillIt2021, rajiAIEverythingWhole2021a}. Standard benchmarks and narrowly-defined tasks are now quickly becoming saturated~\cite{kielaDynabenchRethinkingBenchmarking2021d} and attention is shifting towards testing general-purpose abilities of LLMs~\cite{wangSuperGLUEStickierBenchmark2019,singhGlobalMMLUUnderstanding2025}. The interpretation of such evaluations has become contested, with disagreements about whether results show signs of intelligence~\cite{bubeckSparksArtificialGeneral2023, mccoyEmbersAutoregressionUnderstanding2023b} or emergent abilities~\cite{schaefferAreEmergentAbilities2023a,  weiEmergentAbilitiesLarge2022}, making assessing construct validity all the more crucial.